
de \(C_s\), qui résulte de la connexité géométrique de \(C_t\).

De l'homomorphisme (7.12.1) on déduit grâce à 7.4 un homomorphisme
\begin{equation}\tag{7.12.2}
  \Gr_\bullet^{\mathrm{alg}}(X_t)\longrightarrow \Gr_\bullet^{\mathrm{alg}}(X_s).
\end{equation}

Lorsqu'on suppose de plus \(X_s\) et \(X_t\) réguliers, alors (7.12.1) peut encore s'interpréter comme un homomorphisme d'anneaux
\begin{equation}\tag{7.12.3}
  K^\bullet_{\mathrm{alg}}(X_t)\longrightarrow K^\bullet_{\mathrm{alg}}(X_s),
\end{equation}
qui, lorsque \(X\) est de plus séparé, est un homomorphisme de \(\lambda\)-anneaux, et définit par suite un homomorphisme
\begin{equation}\tag{7.12.4}
  \Gr^\bullet_{\mathrm{alg}}(X_t)\longrightarrow \Gr^\bullet_{\mathrm{alg}}(X_s).
\end{equation}

\textbf{7.12.5.} Soit \(T\) un schéma algébrique lisse et propre sur un corps \(k\). On vérifie facilement qu'un élément \(x\) de \(K^\bullet(T)\), tel qu'il existe un entier \(n\geq1\) avec \(nx\) algébriquement équivalent à zéro, est numériquement équivalent à zéro. Il est plausible que la réciproque soit toujours vraie, du moins si \(T\) est projectif. Dans cette question, on peut manifestement supposer \(k\) algébriquement clos, et remplacer \(K^\bullet(T)\) par \(\Gr^\bullet(T)\) ou \(\Gr_\bullet(T)\), sans modifier la conjecture. On trouve également une conjecture équivalente en travaillant avec l'anneau de Chow \(A(T)\) (cf. XIV 4.1), et on trouve la conjecture plus ou moins classique suivante : si \(Z\) est un cycle algébrique numériquement équivalent à zéro sur \(T\), il existe un entier \(n\geq1\) avec \(nZ\) algébriquement équivalent à zéro. Dans le cas où cette conjecture est correcte pour \(X_t\), il résulterait du résultat de compatibilité de l'homomorphisme (7.3.1) avec l'équivalence algébrique que cet homomorphisme de spécialisation (dans le cas où \(f:X\to S\) est un morphisme projectif et lisse) est aussi compatible avec l'équivalence numérique, comme conjecturé dans 7.9.

